\documentclass{article}
\usepackage{amsmath,amssymb, amsthm}
\usepackage[utf8]{inputenc}
\usepackage[polish]{babel}
\usepackage[T1]{fontenc}
\newtheorem{twierdzenie}{Twierdzenie}

\begin{document}

\section{Czas Uniksowy}
Czas Uniksowy (z ang. \emph{Unix Time}), to sposób liczenia czasu jako liczby sekund od 1 stycznia 1970 r., tak 
zwanej ery uniksa (z ang. \emph{Unix Era}). Jest to system przydatny szczególnie w informatyce, gdzie ułatwia on np. 
prównywanie dat. 

Liczba sekund jest jednak trzymana jedynie na 32 bitach, 
więc największą liczbę jaką może przechować jest 2 147 483 647. Po tylu sekundach od 1 stycznia 1970, czyli 19 stycznia 2038 roku o 03:14:07 UTC, licznik zwyczajnie nie będzie mógł
kontynuować (albo się wyzeruje i wedle komputerów znów będziemy w roku 1970). Tak więc przed upłynięciem tego czasu
trzeba będzie znaleźć rozwiązanie tego problemu. Możemy np. 
zmienić miejsce na 64 bity, wówczas następny taki problem 
pojawi się 9 listopada 584 554 051 223 roku o godzinie 07:00:15 UTC. Oczywiście takie rozwiązanie samo w sobie jest 
problematyczne i tylko teoretyczne.

\section{Matematyka}
\begin{twierdzenie}
    \[\sum_{i=0}^{n}2^i+1=2^{n+1} \]
\end{twierdzenie}
\begin{proof}
    Baza indukcji, dla \(n=0\):

    \[
    \sum_{i=0}^{0}2^0+1 = 2 = 2^1
    \]

    Krok indukcyjny, załóżmy prawdziwość twierdzenia dla n
    i udowodnijmy prawdziwość dla n+1:
    \[
    \sum_{i=0}^{n+1}2^i+1=\sum_{i=0}^{n}2^i+2^{n+1}+1=
    2^{n+1}-1+2^{n+1}+1=2^{n+2}
    \]

\end{proof}

\section{Zwierzątko}

\begin{verbatim}
                    __
         .,-;-;-,. /'_\
       _/_/_/_|_\_\) /
     '-<_><_><_><_>=/\
       `/_/====/_/-'\_\
        ""     ""    ""

\end{verbatim}




\end{document}
